%------------------------------------------------------------------------------%
\section{Zeros de Função}

Os valores de $x$ que satisfazem a equação $f(x)=0$ são chamados
\emph{zeros}\index{Zeros}
de $f$. Esses valores correspondem aos interceptos-$x$ do gráfico da função.

%------------------------------------------------------------------------------%
\begin{example}
  Encontre os zeros da função \(f(x) = x^3+2x^2-5x\)

  \solution
  Para encontrarmos os zeros de $f$ precisamos resolver a equação
  \[
    x^3 + 2x^2 - 5x = 0
    \quad\Rightarrow\quad
    x( x^2 + 2x - 5 ) = 0,
  \]
  portanto
  \[
    x = 0
    \qquad\text{ ou }\qquad
    x^2 + 2x - 5 = 0.
  \]
  Para encontrarmos as raízes de
  \[
    x^2 + 2x- 5 =0
  \]
  usamos a fórmula de Bhaskara, uma vez que se trata de uma equação
  do segundo grau.
  Com isso,
  \[
    x^2 + 2x - 5 = 0
    \quad\Rightarrow\quad
    x = -1+\sqrt{6}
  \]
  ou
  \[
    x = -1-\sqrt{6}.
  \]
  Logo, $f$ possui três zeros, a saber
  \[
    x_1 = 0,             \qquad
    x_2 =-1+\sqrt{6}     \qquad
    \text{ e }           \qquad
    x_3 =-1-\sqrt{6}.
  \]
\end{example}
%------------------------------------------------------------------------------%

%------------------------------------------------------------------------------%
\begin{example}
  Encontre os zeros da função \(g(x) = \frac{x^2-4}{x+3}\)

  \solution
  Para encontrarmos os zeros de $g$ precisamos resolver a equação
  \begin{align*}
    \frac{x^2-4}{x+3} & = 0 \\[2mm]
    x^2-4             & = 0 \\
    x=2               & \;\text{ ou }\; x=-2.
  \end{align*}
  Logo $g$ possui dois zeros, a saber $x_1=-2$ e $x_1=2$.
\end{example}
%------------------------------------------------------------------------------%

%------------------------------------------------------------------------------%